\documentclass[11pt]{article}

% ---------------------------------------------------------------------------
% Packages
% ---------------------------------------------------------------------------
\usepackage[margin=1.1in]{geometry}
\usepackage{lmodern}
\usepackage{amsmath,amssymb,amsthm,mathtools}
\usepackage{booktabs}
\usepackage{enumitem}
\usepackage{xcolor}
\usepackage{graphicx}
\usepackage{caption}
\usepackage[expansion=false]{microtype}
\usepackage[colorlinks=true,linkcolor=blue!55!black,citecolor=blue!50!black,urlcolor=blue!55!black]{hyperref}

% ---------------------------------------------------------------------------
% Theorem environments
% ---------------------------------------------------------------------------
\theoremstyle{plain}
\newtheorem{theorem}{Theorem}[section]
\newtheorem{proposition}[theorem]{Proposition}
\newtheorem{lemma}[theorem]{Lemma}
\newtheorem{corollary}[theorem]{Corollary}
\theoremstyle{definition}
\newtheorem{definition}[theorem]{Definition}
\newtheorem{assumption}{Assumption}[section]
\renewcommand{\theassumption}{A\arabic{assumption}}
\theoremstyle{remark}
\newtheorem*{remark}{Remark}

% ---------------------------------------------------------------------------
% Math macros
% ---------------------------------------------------------------------------
\newcommand{\R}{\mathbb{R}}\newcommand{\N}{\mathbb{N}}\newcommand{\E}{\mathbb{E}}
\newcommand{\Prob}{\mathbb{P}}\newcommand{\one}{\mathbf{1}}\newcommand{\dd}{\,\mathrm{d}}
\newcommand{\Hist}{\mathcal{H}}\newcommand{\Filt}{\mathcal{F}}\newcommand{\conv}{\ast}
\newcommand{\Lap}{\mathcal{L}}
\DeclareMathOperator{\KL}{KL}\DeclareMathOperator*{\argmin}{arg\,min}
\DeclareMathOperator{\Var}{Var}
\newcommand{\ind}[1]{\mathbf{1}\!\left[#1\right]}
\newcommand{\code}[1]{\texttt{#1}}

% ---------------------------------------------------------------------------
\title{\bfseries Calibrating Self-Exciting Investment Dynamics from
Interval-Censored Counts\\[4pt]
\large The Mean Behavior Poisson Process, a Covariate Augmentation,\\
and a Precise Account of Its Limits}
\author{Alberto Gennaro}
\date{\today}

\begin{document}
\maketitle

\begin{abstract}
\noindent
We study the problem of fitting Hawkes (self-exciting) point processes when the
data are \emph{interval-censored}: we observe only the number of events in each
of a sequence of time windows, never the event times themselves. Aggregated
investment activity --- deals closed per week in a sector --- is exactly of this
form. The Hawkes likelihood cannot be evaluated without event times, and the
process lacks independent increments, so the Poisson likelihood does not apply
either. Following Rizoiu et al.\ (2022), we resolve this through the \emph{Mean
Behavior Poisson Process} (MBPP), an inhomogeneous Poisson process whose
deterministic intensity equals the expected Hawkes intensity. We first restate
and \emph{prove in full} the foundational facts: that the expected intensity
solves a Volterra equation of the second kind, that this equation defines a
causal linear time-invariant system with an explicit impulse response, and that
the resulting interval likelihood is a Poisson (equivalently, a Kullback--Leibler
/ Bregman) objective. We then state and prove our central methodological claim:
augmenting the model with covariates --- log-linearly in the baseline
\emph{and} in the excitation --- preserves every structural property the method
relies upon, with one sharp exception that we identify precisely. Baseline
covariates leave the convolution structure intact; excitation covariates replace
the convolution by a general linear Volterra kernel, which remains exactly
solvable by a linear time-varying ODE, but \emph{lose} the transform-based
machinery. The one place the mathematics genuinely breaks is \emph{nonlinear}
self-inhibition, and we explain exactly why. Finally we validate the theory on a
synthetic investment market in two regimes: a deliberately misspecified one
(heavy-tailed power-law triggering, nonlinear self-inhibition, cross-sector
coupling) and a correctly specified one. Under correct specification the
augmented estimator recovers every parameter (branching ratio within $0.03$,
covariate coefficients within $0.08$); under heavy misspecification it still
recovers the sign and significance of the covariate effect and improves held-out
predictive likelihood, while the magnitude is provably inflated --- a bias we
trace to the kernel mismatch.
\end{abstract}

\tableofcontents
\bigskip

% ===========================================================================
\section{Introduction}
\label{sec:intro}

A venture or private-equity data feed rarely gives clean timestamps. What one
reliably has is \emph{counts}: how many financing rounds closed in a sector last
week, how many follow-on investments a fund made last month. The underlying
dynamics are strongly self-exciting --- a marquee round advertises a sector and
pulls in capital, exciting further deals in \emph{similar} sectors --- but with a
twist that distinguishes capital allocation from, say, earthquakes or retweets:
having \emph{just} closed a round, a given firm or fund is briefly
\emph{less} likely to close another. Excitation across similar products coexists
with \emph{self-inhibition}. On top of this, the propensity to transact rises and
falls with exogenous, observable covariates: a sector's popularity, the macro
regime, quiet periods after earnings.

The Hawkes process is the natural model for the self-exciting part. But its
log-likelihood,
\[
\ell(\theta) = \sum_{k}\log\lambda(t_k) - \int_0^T\lambda(s)\dd s,
\]
is a function of the event times $t_k$, which interval censoring destroys. Nor can
we fall back on the Poisson likelihood, because a Hawkes process does \emph{not}
have independent increments: the count in one window is correlated with the count
in the next through the triggering kernel. This is the gap the Mean Behavior
Poisson Process closes.

\paragraph{Contributions and structure.}
This paper is written to be read by a skeptic: every object is defined precisely
and every claim is proved.
\begin{itemize}[leftmargin=1.4em,itemsep=2pt]
\item Section~\ref{sec:prelim} fixes notation and states the standing assumptions
and the two classical facts we build on (the martingale form of the compensator;
Watanabe's characterization of the Poisson process).
\item Section~\ref{sec:mbpp} defines the MBPP, proves that the expected intensity
solves the Volterra equation $\xi=s+\xi\conv\varphi$ (Lemma~\ref{lem:mean}) and
that the expected compensator solves the matching equation (Lemma~\ref{lem:comp}),
and shows the MBPP is genuinely Poisson (Proposition~\ref{prop:poisson}).
\item Section~\ref{sec:solve} solves the equation two ways: by Laplace transform
(Theorem~\ref{thm:laplace}) and by the impulse response
$E=\delta+\sum_{n\ge1}\varphi^{\conv n}$ (Theorem~\ref{thm:impulse}), with the
convergence proof in full, and specializes to the exponential kernel.
\item Section~\ref{sec:estim} derives the interval-censored likelihood
(Proposition~\ref{prop:icll}) and its Bregman/KL reading
(Proposition~\ref{prop:bregman}).
\item Section~\ref{sec:augment} is the methodological core: we prove that baseline
covariates preserve the construction \emph{exactly}
(Theorem~\ref{thm:baseline}), that excitation covariates preserve it up to the
loss of convolution but remain exactly solvable
(Theorem~\ref{thm:excitation}), and we delimit precisely the one case ---
nonlinear inhibition --- where the mean-field identity fails
(Proposition~\ref{prop:nonlinear}).
\item Section~\ref{sec:experiments} validates everything on a synthetic
investment market, misspecified and correctly specified.
\end{itemize}
Sections~\ref{sec:prelim}--\ref{sec:estim} follow Rizoiu et al.\ (2022); the
proofs are reproduced so the document is self-contained. Section~\ref{sec:augment}
and the experiments are our contribution.

% ===========================================================================
\section{Preliminaries and standing assumptions}
\label{sec:prelim}

\subsection{Counting processes, intensity, compensator}

Let $(\Omega,\Filt,(\Filt_t)_{t\ge0},\Prob)$ be a filtered probability space
satisfying the usual conditions. A simple point process on $[0,\infty)$ is
identified with its counting process $N=(N(t))_{t\ge0}$, $N(t)=\#\{k:t_k\le t\}$,
right-continuous, with unit jumps. We write $\Hist_t=\{t_k:t_k<t\}$ for the
strict history.

\begin{definition}[conditional intensity]
A nonnegative $(\Filt_t)$-predictable process $\lambda(t)$ is the
\emph{conditional intensity} of $N$ if
$\Lambda(t)=\int_0^t\lambda(s)\dd s$ is the $(\Filt_t)$-compensator of $N$; that
is, $M(t)=N(t)-\Lambda(t)$ is an $(\Filt_t)$-local martingale. Heuristically
$\E[\dd N(t)\mid\Filt_{t^-}]=\lambda(t)\dd t$.
\end{definition}

We will use repeatedly the following elementary but crucial consequence.

\begin{lemma}[expectation of the counting measure]
\label{lem:campbell}
If $\lambda$ is the intensity of $N$ and $\E[\Lambda(t)]<\infty$ for all $t$, then
for every bounded measurable $g$,
\[
\E\!\left[\int_0^t g(s)\dd N(s)\right]=\E\!\left[\int_0^t g(s)\lambda(s)\dd s\right].
\]
In particular $\E[\dd N(s)]=\E[\lambda(s)]\dd s$ as measures on $[0,t]$.
\end{lemma}
\begin{proof}
For $g=\ind{(a,b]}$ the statement is $\E[N(b)-N(a)]=\E[\Lambda(b)-\Lambda(a)]$,
which is the (optional sampling) martingale property of $M=N-\Lambda$ together
with $\E[\Lambda(t)]<\infty$. Linearity extends it to simple $g$, and bounded
convergence to bounded measurable $g$.
\end{proof}

\subsection{The Poisson process and Watanabe's theorem}

\begin{definition}[inhomogeneous Poisson process]
Given a locally integrable $\xi:[0,\infty)\to[0,\infty)$ with
$\Xi(t)=\int_0^t\xi$, an inhomogeneous Poisson process of intensity $\xi$ is a
counting process with independent increments such that
$N(b)-N(a)\sim\mathrm{Poisson}\!\big(\Xi(b)-\Xi(a)\big)$ for all $a<b$.
\end{definition}

The following classical characterization is the engine of the whole approach: it
says that a simple point process whose compensator is \emph{deterministic} is
necessarily Poisson. We state it precisely because everything downstream depends
on it.

\begin{theorem}[Watanabe]
\label{thm:watanabe}
Let $N$ be a simple point process with $(\Filt_t)$-compensator $\Lambda$. If
$\Lambda$ is continuous and \emph{deterministic} --- $\Lambda(t)=\Xi(t)$ a.s.\ for
a deterministic continuous nondecreasing $\Xi$ --- then $N$ is an inhomogeneous
Poisson process with mean function $\Xi$; i.e.\ increments are independent and
$N(b)-N(a)\sim\mathrm{Poisson}(\Xi(b)-\Xi(a))$.
\end{theorem}
\begin{proof}[Proof (sketch, standard)]
For fixed $u\in\R$ and $a<b$ consider $\phi(t)=\exp\{iu(N(t)-N(a))\}$ on $t\in[a,b]$.
The jump of $\phi$ at a point of $N$ multiplies it by $e^{iu}-1$; writing the
Doob--Meyer decomposition of the finite-variation pure-jump process $\phi$ and
using that its compensator is $\int(e^{iu}-1)\phi(s^-)\dd\Xi(s)$ (here $\Xi$
deterministic), one gets the ODE
$\E[\phi(b)]=1+(e^{iu}-1)\int_a^b\E[\phi(s)]\dd\Xi(s)$, whose solution is
$\E[e^{iu(N(b)-N(a))}]=\exp\{(e^{iu}-1)(\Xi(b)-\Xi(a))\}$. This is the
characteristic function of $\mathrm{Poisson}(\Xi(b)-\Xi(a))$. Determinism of
$\Xi$ makes the increment characteristic functions factorize over disjoint
intervals, giving independence.
\end{proof}

\subsection{The Hawkes process}

\begin{definition}[linear Hawkes process]
\label{def:hawkes}
A linear Hawkes process is a simple point process whose intensity is
\begin{equation}
\lambda(t)=s(t)+\sum_{t_k<t}\varphi(t-t_k)=s(t)+\int_0^t\varphi(t-u)\dd N(u),
\label{eq:hawkes}
\end{equation}
with \emph{exogenous} (immigrant) rate $s(t)\ge0$ and \emph{triggering kernel}
$\varphi:[0,\infty)\to[0,\infty)$, $\varphi(t)=0$ for $t<0$.
\end{definition}

We impose the following throughout Sections~\ref{sec:mbpp}--\ref{sec:estim}.

\begin{assumption}[nonnegativity]\label{ass:nonneg}
$s(t)\ge0$ and $\varphi(t)\ge0$; the intensity \eqref{eq:hawkes} is thus
automatically nonnegative and \eqref{eq:hawkes} is well posed (linear, no
truncation).
\end{assumption}
\begin{assumption}[subcriticality]\label{ass:subcrit}
The \emph{branching ratio} $\kappa:=\int_0^\infty\varphi(t)\dd t$ satisfies
$\kappa<1$.
\end{assumption}
\begin{assumption}[local integrability]\label{ass:integ}
$s$ is locally integrable and $\sup_{t\le T}\E[\lambda(t)]<\infty$ for each $T$.
\end{assumption}

Assumption~\ref{ass:nonneg} is what makes the process \emph{linear}: the intensity
is an affine functional of the past with no rectification. Assumption~\ref{ass:subcrit}
is the stability condition --- under the cluster (Hawkes--Oakes) representation,
each event spawns on average $\kappa$ direct offspring, so $\kappa<1$ makes total
progeny a.s.\ finite and the process non-explosive; it is also exactly the
condition under which the Neumann series of Section~\ref{sec:solve} converges.
We flag now, because it is the hinge of Section~\ref{sec:augment}, that
Assumption~\ref{ass:nonneg} is what will fail for self-inhibition.

% ===========================================================================
\section{The Mean Behavior Poisson Process}
\label{sec:mbpp}

\subsection{Definition and the mean-intensity equation}

The idea is to replace the random Hawkes intensity by its expectation and ask what
process \emph{has} that expectation as a genuine (deterministic) intensity.

\begin{definition}[MBPP]
\label{def:mbpp}
The \emph{Mean Behavior Poisson Process} associated with the Hawkes process
\eqref{eq:hawkes} is the inhomogeneous Poisson process whose deterministic
intensity is the expected Hawkes intensity,
\begin{equation}
\xi(t):=\E_{\Hist_t}[\lambda(t)].
\label{eq:xidef}
\end{equation}
\end{definition}

The content of the definition is that $\xi$ obeys a closed equation.

\begin{lemma}[mean-intensity Volterra equation]
\label{lem:mean}
Under Assumptions~\ref{ass:nonneg}--\ref{ass:integ}, $\xi$ in \eqref{eq:xidef}
satisfies
\begin{equation}
\xi(t)=s(t)+\int_0^t\varphi(t-u)\,\xi(u)\dd u = s(t)+(\xi\conv\varphi)(t).
\label{eq:mbpp}
\end{equation}
\end{lemma}
\begin{proof}
Take expectations in \eqref{eq:hawkes}. The exogenous term is deterministic. For
the endogenous term, by Lemma~\ref{lem:campbell} with $g_t(u)=\varphi(t-u)$
(bounded on $[0,t]$ once we note $\varphi$ is locally bounded; the general case
follows by monotone approximation since $\varphi\ge0$),
\[
\E\!\left[\int_0^t\varphi(t-u)\dd N(u)\right]
=\int_0^t\varphi(t-u)\,\E[\lambda(u)]\dd u
=\int_0^t\varphi(t-u)\,\xi(u)\dd u .
\]
Adding the two terms gives \eqref{eq:mbpp}. The interchange of expectation and
integral is justified by Tonelli's theorem ($\varphi\ge0$, $\E[\lambda]\ge0$) and
finiteness from Assumption~\ref{ass:integ}.
\end{proof}

Equation~\eqref{eq:mbpp} is a \emph{Volterra integral equation of the second
kind} with convolution kernel. Note what the proof used: \emph{only} linearity of
\eqref{eq:hawkes} in the past and the compensator identity of
Lemma~\ref{lem:campbell}. It did \emph{not} use the specific form of $s$, nor
that the kernel is a convolution. We will cash this observation in
Section~\ref{sec:augment}.

\subsection{The compensator obeys the same equation}

\begin{lemma}[compensator Volterra equation]
\label{lem:comp}
Let $\Xi(t):=\E_{\Hist_t}[\Lambda(t)]=\int_0^t\xi$ and $S(t):=\int_0^t s$. Then
\begin{equation}
\Xi(t)=S(t)+\int_0^t\varphi(t-u)\,\Xi(u)\dd u=S(t)+(\Xi\conv\varphi)(t).
\label{eq:mbppcomp}
\end{equation}
\end{lemma}
\begin{proof}
Since $\lambda$ is bounded in expectation we may pass the expectation through the
time integral (Fubini, justified by Assumption~\ref{ass:integ}):
$\Xi(t)=\int_0^t\E[\lambda(z)]\dd z=\int_0^t\xi(z)\dd z$. Substitute
\eqref{eq:mbpp}:
\[
\Xi(t)=\int_0^t s(z)\dd z+\int_0^t\!\!\int_0^z\xi(z-y)\varphi(y)\dd y\dd z .
\]
In the double integral, reverse the order over the region $\{0\le y\le z\le t\}$:
\[
\int_0^t\!\!\int_0^z\xi(z-y)\varphi(y)\dd y\dd z
=\int_0^t\varphi(y)\!\int_y^t\xi(z-y)\dd z\,\dd y
=\int_0^t\varphi(y)\!\int_0^{t-y}\xi(x)\dd x\,\dd y,
\]
where we substituted $x=z-y$. The inner integral is $\Xi(t-y)$, giving
$\Xi(t)=S(t)+\int_0^t\varphi(y)\Xi(t-y)\dd y$, which is \eqref{eq:mbppcomp}.
\end{proof}

That $\xi$ and $\Xi$ solve the \emph{same} convolution equation (with forcing $s$
resp.\ $S$) is what lets a single solution method serve both fitting and
forecasting.

\subsection{The MBPP is a Poisson process}

\begin{proposition}
\label{prop:poisson}
Suppose \eqref{eq:mbpp} has a unique locally integrable nonnegative solution
$\xi$ (Theorem~\ref{thm:impulse} below gives it explicitly under
Assumption~\ref{ass:subcrit}). Then the process of Definition~\ref{def:mbpp} is
well defined, and it is an inhomogeneous Poisson process with mean function
$\Xi=\int_0^\cdot\xi$. Consequently its increments are independent and
$N(o_{i-1},o_i]\sim\mathrm{Poisson}\big(\Xi(o_i)-\Xi(o_{i-1})\big)$.
\end{proposition}
\begin{proof}
By construction the MBPP has the deterministic continuous compensator $\Xi$.
Apply Watanabe's Theorem~\ref{thm:watanabe}.
\end{proof}

This is the decisive structural gain. The Hawkes process has a \emph{random}
compensator (it depends on the realized history) and hence correlated, non-Poisson
increments. Its mean-behavior companion has a \emph{deterministic} compensator and
therefore, by Watanabe, \emph{independent} Poisson increments --- precisely the
property interval censoring needs.

\begin{remark}[one-to-one correspondence]
The pair $(s,\varphi)$ determines both $\lambda$ (the Hawkes intensity) and $\xi$
(the MBPP intensity). We therefore identify a Hawkes process with its MBPP and
fit the latter; recovering $(s,\varphi)$ recovers the Hawkes parameters.
\end{remark}

% ===========================================================================
\section{Solving the mean-behavior equation}
\label{sec:solve}

Equation~\eqref{eq:mbpp} admits a closed solution only in special cases. We give
the two methods of Rizoiu et al.\ (2022): a transform method, exact but
restrictive, and an impulse-response method, general.

\subsection{Laplace-transform solution}

\begin{theorem}[transform solution]
\label{thm:laplace}
If the Laplace transforms $\Lap\{s\}$ and $\Lap\{\varphi\}$ exist and
$\Lap\{\varphi\}(\omega)\ne1$ on the region of convergence, then
\begin{equation}
\xi(t)=\Lap^{-1}\!\left\{\frac{\Lap\{s\}(\omega)}{1-\Lap\{\varphi\}(\omega)}\right\}(t)
\label{eq:laplace}
\end{equation}
solves \eqref{eq:mbpp}, whenever the inverse transform exists.
\end{theorem}
\begin{proof}
The Laplace transform turns convolution into product: $\Lap\{\xi\conv\varphi\}
=\Lap\{\xi\}\Lap\{\varphi\}$. Transforming \eqref{eq:mbpp},
$\Lap\{\xi\}=\Lap\{s\}+\Lap\{\xi\}\Lap\{\varphi\}$, so
$\Lap\{\xi\}(1-\Lap\{\varphi\})=\Lap\{s\}$, and \eqref{eq:laplace} follows by
inverting. Existence of the inverse is exactly the stated hypothesis.
\end{proof}

The method is exact but fragile: the immigrant functions we will need involve
Dirac masses (no classical transform), and power-law or Rayleigh kernels have no
elementary inverse transform. This motivates the second method.

\subsection{Impulse-response solution}

We first record that \eqref{eq:mbpp} defines a causal linear time-invariant (LTI)
system mapping input $s$ to output $\xi$, then compute its impulse response.

\begin{theorem}[impulse response]
\label{thm:impulse}
Under Assumption~\ref{ass:subcrit}, the solution operator $s\mapsto\xi$ of
\eqref{eq:mbpp} is causal, linear and time-invariant, and is characterized by its
impulse response
\begin{equation}
E(t)=\delta(t)+h(t),\qquad h(t)=\sum_{n=1}^\infty\varphi^{\conv n}(t),
\label{eq:impulse}
\end{equation}
where $\varphi^{\conv n}$ is the $n$-fold convolution of $\varphi$. The series
converges in $L^1$.
\end{theorem}
\begin{proof}
\emph{Linearity} is immediate: if $\xi_i$ solves \eqref{eq:mbpp} with forcing
$s_i$, then $a_1\xi_1+a_2\xi_2$ solves it with forcing $a_1s_1+a_2s_2$.
\emph{Time-invariance}: let $s\mapsto\xi$ and consider the shifted input
$s(\cdot-t_0)$. With $t'=t-t_0$,
\[
\xi(t-t_0)=s(t-t_0)+\int_0^t\xi(\tau)\varphi(t-t_0-\tau)\dd\tau
=s(t')+\int_0^{t'+t_0}\xi(\tau)\varphi(t'-\tau)\dd\tau,
\]
and splitting the integral at $t'$, the part over $(t',t'+t_0)$ vanishes because
$\varphi(t'-\tau)=0$ for $\tau>t'$ (causality of $\varphi$); hence the response to
$s(\cdot-t_0)$ is $\xi(\cdot-t_0)$. \emph{Causality} of $\xi$ follows from that of
$\varphi$ and $s$.

For the impulse response, feed the unit impulse $\delta$ as forcing; the output
$E$ satisfies $E=\delta+E\conv\varphi$. Writing $E=\delta+h$ and using
$\delta\conv\varphi=\varphi$,
\[
\delta+h=\delta+(\delta+h)\conv\varphi=\delta+\varphi+h\conv\varphi
\;\Longrightarrow\; h=\varphi+h\conv\varphi.
\]
Iterating, $h=\varphi+\varphi^{\conv2}+\dots+\varphi^{\conv n}+h\conv\varphi^{\conv n}$,
so $h=\sum_{n\ge1}\varphi^{\conv n}$ provided the remainder vanishes. By Young's
convolution inequality applied to the $L^1$ norm $\|f\|_1=\int_0^\infty|f|$,
\[
\int_0^\infty\varphi^{\conv n}(t)\dd t\le\Big(\int_0^\infty\varphi(t)\dd t\Big)^n=\kappa^n\xrightarrow[n\to\infty]{}0,
\]
using $\kappa<1$ (Assumption~\ref{ass:subcrit}); since $\varphi^{\conv n}\ge0$,
$\|\varphi^{\conv n}\|_1\to0$, so the series converges in $L^1$ and the remainder
$h\conv\varphi^{\conv n}\to0$.
\end{proof}

\begin{corollary}[solution for arbitrary forcing]
\label{cor:solution}
For any (generalized) forcing $s$, $\xi=E\conv s=s+h\conv s$ solves
\eqref{eq:mbpp}.
\end{corollary}
\begin{proof}
$\xi=E\conv s=(\delta+h)\conv s=s+h\conv s$. Using $h=E\conv\varphi$ (from
$E=\delta+E\conv\varphi$), $\xi=s+(E\conv s)\conv\varphi=s+\xi\conv\varphi$.
\end{proof}

The impulse response handles Dirac immigrants and any kernel; it requires no
transform inversion and reuses $E$ across many forcings $s$ --- the property we
exploit when refitting under changing covariates.

\subsection{The exponential kernel, in closed form}
\label{sec:exp}

Take the exponential kernel
$\varphi(t)=\kappa\theta e^{-\theta t}\ind{t>0}$, so
$\int_0^\infty\varphi=\kappa$. Its $n$-fold convolution is, by induction,
\[
\varphi^{\conv n}(t)=\kappa^n\theta^n e^{-\theta t}\frac{t^{n-1}}{(n-1)!},\qquad t\ge0,
\]
(the base case is immediate; the inductive step is a Gamma integral). Summing,
\begin{equation}
h(t)=\sum_{n\ge1}\varphi^{\conv n}(t)
=\kappa\theta e^{-\theta t}\sum_{n\ge0}\frac{(\kappa\theta t)^n}{n!}
=\kappa\theta e^{(\kappa-1)\theta t},\qquad t\ge0,
\label{eq:hexp}
\end{equation}
so the impulse response is $E(t)=\delta(t)+\kappa\theta e^{(\kappa-1)\theta t}
\ind{t>0}$. For a constant immigrant rate $s\equiv\mu$ one gets, by
Corollary~\ref{cor:solution},
\[
\xi(t)=\mu+\int_0^t h(t-u)\,\mu\dd u
=\mu\Big(1+\tfrac{\kappa}{1-\kappa}\big(1-e^{-(1-\kappa)\theta t}\big)\Big)
\xrightarrow[t\to\infty]{}\frac{\mu}{1-\kappa},
\]
recovering the stationary mean rate $\mu/(1-\kappa)$. We adopt the
$(\kappa,\theta)$ parametrization throughout: $\kappa\in(0,1)$ is the branching
ratio (well identified from counts) and $\theta>0$ the decay (weakly identified;
see Section~\ref{sec:estim}). The corresponding $\alpha=\kappa\theta$,
$\beta=\theta$ recover the usual $\alpha e^{-\beta t}$ form.

\begin{corollary}[compensator over an interval]
\label{cor:compint}
With $E=\delta+h$ the impulse response and $S=\int_0^\cdot s$, the MBPP
compensator over $(x,y]$ is
\begin{equation}
\Xi(x,y]=(E\conv S)(y)-(E\conv S)(x).
\label{eq:compint}
\end{equation}
\end{corollary}
\begin{proof}
By Lemma~\ref{lem:comp} the compensator obeys the same convolution equation as
$\xi$ with forcing $S$; apply Corollary~\ref{cor:solution} to get
$\Xi=E\conv S$, then difference at the endpoints.
\end{proof}

% ===========================================================================
\section{Estimation from interval-censored data}
\label{sec:estim}

\subsection{The interval-censored likelihood}

We observe ordered endpoints $0=o_0<o_1<\dots<o_m=T$ and the \emph{counts}
$C_i:=N(o_{i-1},o_i]$, and nothing else. Because the MBPP is Poisson
(Proposition~\ref{prop:poisson}), the counts are independent with
$C_i\sim\mathrm{Poisson}(\Xi_i)$, $\Xi_i:=\Xi(o_{i-1},o_i;\theta]$.

\begin{proposition}[interval-censored log-likelihood]
\label{prop:icll}
The negative log-likelihood of the counts under the MBPP is, up to an additive
constant independent of $\theta$,
\begin{equation}
\mathcal L_{\mathrm{IC}}(\theta)=\sum_{i=1}^m\Xi_i(\theta)-\sum_{i=1}^m C_i\log\Xi_i(\theta).
\label{eq:icll}
\end{equation}
\end{proposition}
\begin{proof}
By independence and the Poisson pmf,
$\Prob(C_1,\dots,C_m)=\prod_i e^{-\Xi_i}\Xi_i^{C_i}/C_i!$. Taking $-\log$,
$\mathcal L=\sum_i(\Xi_i-C_i\log\Xi_i+\log C_i!)$; the term $\sum_i\log C_i!$
depends only on the data and is dropped under minimization over $\theta$.
\end{proof}

\subsection{The Bregman / Kullback--Leibler reading}

Write $C=(C_1,\dots,C_m)$, $\Xi(\theta)=(\Xi_1,\dots,\Xi_m)$. Recall the Bregman
divergence of a strictly convex generator $\phi$,
$B_\phi(x,y)=\phi(x)-\phi(y)-\langle x-y,\nabla\phi(y)\rangle\ge0$.

\begin{proposition}[IC-LL is a KL divergence]
\label{prop:bregman}
With the generator $\phi(x)=\sum_i(x_i\log x_i-x_i)$ one has, up to a
$\theta$-independent constant,
\begin{equation}
\mathcal L_{\mathrm{IC}}(\theta)=\KL\big(C,\Xi(\theta)\big)+\mathrm{const},
\qquad
\argmin_\theta\mathcal L_{\mathrm{IC}}=\argmin_\theta\KL\big(C,\Xi(\theta)\big).
\label{eq:kl}
\end{equation}
Replacing $\phi$ by the squared generator $\tfrac12\|x\|^2$ yields instead the
least-squares loss $\sum_i(C_i-\Xi_i)^2$.
\end{proposition}
\begin{proof}
For $\phi(x)=\sum_i(x_i\log x_i-x_i)$, $\nabla\phi(y)_i=\log y_i$, so
\[
B_\phi(C,\Xi)=\sum_i\big(C_i\log C_i-C_i\big)-\big(\Xi_i\log\Xi_i-\Xi_i\big)-\log\Xi_i\,(C_i-\Xi_i)
=\sum_i\big(\Xi_i-C_i\log\Xi_i\big)+\underbrace{\sum_i(C_i\log C_i-C_i)}_{\text{const in }\theta}.
\]
The first sum is \eqref{eq:icll}; hence $\mathcal L_{\mathrm{IC}}=B_\phi(C,\Xi)+\mathrm{const}=\KL(C,\Xi)+\mathrm{const}$.
For $\phi=\tfrac12\|x\|^2$, $\nabla\phi(y)=y$ and $B_\phi(C,\Xi)=\tfrac12\|C-\Xi\|^2$.
\end{proof}

The choice of divergence is the choice of noise model: by the regular
Bregman--exponential-family bijection, KL corresponds to a Poisson observation per
window and least squares to a Gaussian one. KL is the correct default here because
a Poisson window has variance equal to its mean, so \eqref{eq:kl} automatically
weights window $i$ by $1/\Xi_i$; the squared loss assumes homoscedastic Gaussian
noise and is the (corrected) objective underlying the earlier HIP heuristic.

\subsection{History-aware numerical compensator}

When $h$ has no closed form one approximates $\Xi$ numerically. Replacing the
MBPP counting process by its piecewise-constant lower bound $M_D^-$ on a grid
$\{d_j\}$ and using $\Xi(t)=\E[M(t)]$ for a Poisson process gives the convergent
approximation $\Xi(t)\approx S(t)+\sum_j\E[M(d_{j-1},d_j]]\int\varphi$, which
becomes exact as the grid refines. Its practical virtue is that the increments
$\E[M(d_{j-1},d_j]]$ can be replaced by \emph{observed} past counts, turning the
same formula into a one-step-ahead \emph{forecaster}; we use this below only
implicitly, through held-out evaluation.

\subsection{Identifiability: what counts can and cannot tell us}
\label{sec:identif}

\begin{remark}[the $\kappa$--$\theta$ ridge]
\label{rem:ridge}
For the exponential kernel the total mass $\int_0^\infty\varphi=\kappa$ is
independent of $\theta$; changing $\theta$ at fixed $\kappa$ \emph{reshapes}
$\varphi$ in time without changing its mass. A censoring window reports only
$\Xi_i=\int_{o_{i-1}}^{o_i}\xi$, which averages over exactly the within-window
timing that carries $\theta$. Consequently the Fisher information
$\mathcal I_{jk}=\sum_i\Xi_i^{-1}(\partial_j\Xi_i)(\partial_k\Xi_i)$ is
near-degenerate along curves of constant $\kappa$: $\kappa$ (a level effect) is
well identified, $\theta$ (a timing effect) is weakly identified, the more so the
coarser the windows. The covariate coefficients below are level effects and, like
$\kappa$, are well identified; we therefore report $\kappa$, $\gamma$, $\delta$
with confidence and treat $\theta$ as a weak nuisance.
\end{remark}

% ===========================================================================
\section{The covariate augmentation, and why it preserves the mathematics}
\label{sec:augment}

We now add covariates and prove, with the precision a skeptic is owed, exactly
which structural properties survive. There are two augmentations and one
boundary.

\subsection{Baseline covariates: the construction is preserved exactly}

Let $X:[0,\infty)\to\R^p$ be an observable, deterministic covariate path and set
the immigrant rate to the log-linear form
\begin{equation}
s(t)=\exp\!\big(\gamma_0+\gamma^\top X(t)\big)\ge0.
\label{eq:basecov}
\end{equation}

\begin{theorem}[baseline soundness]
\label{thm:baseline}
Under \eqref{eq:basecov} with $X$ deterministic and locally bounded, every result
of Sections~\ref{sec:mbpp}--\ref{sec:estim} holds verbatim: the mean intensity
solves \eqref{eq:mbpp} and the compensator \eqref{eq:mbppcomp} with this $s$; the
impulse response $E=\delta+h$ is \emph{unchanged} (it depends only on $\varphi$);
the solution is $\xi=E\conv s$; the MBPP is Poisson; and the likelihood is
\eqref{eq:icll}. If $X$ is piecewise constant then $s$ is piecewise constant and
$\Xi$ is available in closed form via \eqref{eq:compint}.
\end{theorem}
\begin{proof}
Inspect the proofs of Lemmas~\ref{lem:mean}--\ref{lem:comp}: they use only that
$s$ is a fixed deterministic, nonnegative, locally integrable function --- never
its functional form. Equation~\eqref{eq:basecov} is such a function (a continuous
positive transform of a locally bounded path), so \eqref{eq:mbpp},
\eqref{eq:mbppcomp} hold. The impulse response (Theorem~\ref{thm:impulse}) is
defined by $\varphi$ alone, so it is untouched; $\xi=E\conv s$ by
Corollary~\ref{cor:solution}. The compensator is deterministic, so Watanabe gives
Poisson increments and \eqref{eq:icll}. For piecewise-constant $X$, $s$ is a
step function and the convolution $E\conv S$ in \eqref{eq:compint} is elementary.
\end{proof}

In words: baseline covariates are \emph{free}. They change the forcing $s$ of an
LTI system without touching the system, so the entire toolbox transfers without a
new approximation. This is the interval-censored analogue of a log-linear
intensity baseline in event-time Hawkes models.

\subsection{Excitation covariates: convolution is lost, exact solvability is not}

Now let the \emph{strength} of triggering depend on an observable path
$Z:[0,\infty)\to\R^q$. With an exponential base shape and the receiving-time
convention, each past event contributes, at present time $t$,
\begin{equation}
\text{instantaneous excitation}=\kappa\theta\,\exp\!\big(\delta^\top Z(t)\big)\,e^{-\theta(t-u)},
\label{eq:exccov}
\end{equation}
so the time-varying branching ratio is $\kappa(t)=\kappa\,e^{\delta^\top Z(t)}$.
The Hawkes intensity is still \emph{linear} in the history,
$\lambda(t)=s(t)+\int_0^t K(t,u)\dd N(u)$ with the \emph{separable, non-convolution}
kernel
\begin{equation}
K(t,u)=\kappa\theta\,e^{\delta^\top Z(t)}\,e^{-\theta(t-u)}=m(t)\,g(t-u),\quad
m(t):=\kappa\theta\,e^{\delta^\top Z(t)},\; g(\tau):=e^{-\theta\tau}.
\label{eq:sepkernel}
\end{equation}

\begin{theorem}[excitation soundness]
\label{thm:excitation}
Assume $Z$ is deterministic, locally bounded, and the modulated branching ratio
is subcritical, $\sup_{t\le T}\kappa\,e^{\delta^\top Z(t)}<1$. Then:
\begin{enumerate}[label=\textup{(\roman*)},leftmargin=2.2em,itemsep=2pt]
\item the expected intensity satisfies the linear Volterra equation
\begin{equation}
\xi(t)=s(t)+\int_0^t K(t,u)\,\xi(u)\dd u;
\label{eq:exc-volterra}
\end{equation}
\item \eqref{eq:exc-volterra} has a unique locally bounded solution;
\item the associated MBPP is Poisson, so the likelihood is still \eqref{eq:icll};
\item the convolution structure is \emph{lost} --- $K(t,u)$ is not a function of
$t-u$ alone --- so Theorems~\ref{thm:laplace}--\ref{thm:impulse} no longer apply;
\item nevertheless, for the exponential base shape \eqref{eq:sepkernel} the
equation is equivalent to the scalar linear time-varying ODE
\begin{equation}
y'(t)=s(t)+\theta\big(\kappa\,e^{\delta^\top Z(t)}-1\big)\,y(t),\qquad
\xi(t)=s(t)+\kappa\theta\,e^{\delta^\top Z(t)}\,y(t),
\label{eq:ltv}
\end{equation}
with $y(0)=0$; for piecewise-constant $Z$ this is solved \emph{exactly}, segment
by segment, in closed form.
\end{enumerate}
\end{theorem}
\begin{proof}
(i) The mean-field derivation of Lemma~\ref{lem:mean} used only linearity of the
intensity in the past and Lemma~\ref{lem:campbell}; with the (still
deterministic-given-$Z$, still nonnegative) kernel $K$, the same computation gives
$\xi(t)=s(t)+\int_0^t K(t,u)\E[\lambda(u)]\dd u=s(t)+\int_0^t K(t,u)\xi(u)\dd u$.

(ii) Equation~\eqref{eq:exc-volterra} is a linear Volterra equation of the second
kind with kernel bounded on the triangle $0\le u\le t\le T$. Its resolvent is the
Neumann series $\sum_{n\ge0}K^{\conv n}$ (iterated kernels); a standard bound gives
$|K_n(t,u)|\le \bar m^{\,n}(t-u)^{n-1}/(n-1)!$ with $\bar m=\sup_t m(t)$, whose sum
converges absolutely and uniformly on $[0,T]$. Hence a unique locally bounded
solution exists. (Subcriticality is not needed for existence on bounded $[0,T]$;
it is what keeps the solution bounded as $T\to\infty$.)

(iii) The MBPP is defined to have the deterministic compensator
$\Xi=\int_0^\cdot\xi$; Watanabe (Theorem~\ref{thm:watanabe}) gives independent
Poisson increments, hence \eqref{eq:icll}.

(iv) Immediate from \eqref{eq:sepkernel}: $K(t,u)$ depends on $t$ through $m(t)$,
not only through $t-u$, so it is not a convolution and the transform identity
$\Lap\{\xi\conv\varphi\}=\Lap\{\xi\}\Lap\{\varphi\}$ has no analogue.

(v) Define the state $y(t):=\int_0^t e^{-\theta(t-u)}\xi(u)\dd u$. Then
$\int_0^t K(t,u)\xi(u)\dd u=\kappa\theta e^{\delta^\top Z(t)}y(t)$, so by (i)
$\xi(t)=s(t)+\kappa\theta e^{\delta^\top Z(t)}y(t)$, the second equation of
\eqref{eq:ltv}. Differentiating $y$ (Leibniz),
$y'(t)=\xi(t)-\theta y(t)=s(t)+\theta(\kappa e^{\delta^\top Z(t)}-1)y(t)$, the
first equation. On any interval where $Z$ (hence the coefficient
$r(t)=\theta(1-\kappa e^{\delta^\top Z}))$ is constant, \eqref{eq:ltv} is a
constant-coefficient linear ODE with the exact solution
$y(t)=y(t_0)e^{-r(t-t_0)}+\tfrac{s}{r}(1-e^{-r(t-t_0)})$ (and the obvious limit
when $r=0$), from which the within-segment contribution to $\Xi=\int\xi$ is
computed in closed form. Concatenating segments yields $\Xi$ exactly.
\end{proof}

\begin{remark}[what was gained and lost]
The mean-field identity, the Poisson structure, and the likelihood \eqref{eq:icll}
survive intact, because they rest on linearity and Watanabe, not on convolution.
What is lost is only the \emph{computational} convenience of transforms and a
single reusable impulse response; the price is paid by integrating a
one-dimensional ODE, which for piecewise-constant covariates is itself closed-form.
The estimator maximizing \eqref{eq:icll} over $(\kappa,\theta,\gamma,\delta)$ is
\emph{not} convex in general (the map $\delta\mapsto\Xi$ is nonlinear), so we use
multi-start local optimization; identifiability of $\delta$ requires that the
covariate actually varies while events occur.
\end{remark}

\subsection{The boundary: nonlinear self-inhibition genuinely breaks the identity}

Our investment narrative wants \emph{self-inhibition}: just after a deal, the same
sector is less likely to transact. The faithful way to encode endogenous
self-suppression is a \emph{negative} self-kernel together with a rectifying link
that keeps the intensity nonnegative,
\begin{equation}
\lambda(t)=\Big(s(t)+\textstyle\sum_{t_k<t}\Phi(t-t_k)\Big)^+,\qquad (x)^+=\max(0,x),
\label{eq:nonlinear}
\end{equation}
with $\Phi$ allowed to take negative values (self-inhibition, cross-excitation
positive). This violates Assumption~\ref{ass:nonneg}: the intensity is no longer an
\emph{affine} functional of the past.

\begin{proposition}[the mean-field equation fails under rectification]
\label{prop:nonlinear}
For the nonlinear intensity \eqref{eq:nonlinear} the expected intensity
$\xi=\E[\lambda]$ does \emph{not} in general satisfy a closed linear equation of
the form \eqref{eq:mbpp} or \eqref{eq:exc-volterra}. Quantitatively, writing
$Y(t)=s(t)+\sum_{t_k<t}\Phi(t-t_k)$ for the pre-rectification field,
\[
\xi(t)=\E[(Y(t))^+]\;\ge\;\big(\E[Y(t)]\big)^+,
\]
with strict inequality whenever $Y(t)$ has positive probability of being negative;
in particular the linear surrogate $\E[Y(t)]$ \emph{underestimates} $\xi(t)$.
\end{proposition}
\begin{proof}
$x\mapsto(x)^+$ is convex, so Jensen gives
$\E[(Y)^+]\ge(\E[Y])^+$. The map is strictly convex across $0$, so the inequality
is strict unless $Y(t)\ge0$ a.s.\ or $Y(t)\le 0$ a.s. Because the expectation and
the rectifier do not commute, taking expectations of \eqref{eq:nonlinear} cannot
produce a closed equation in $\xi=\E[\lambda]$ alone: the right-hand side involves
the full law of $Y(t)$, not merely its mean.
\end{proof}

This is the precise limit of the method. As long as the model is \emph{linear} ---
covariates entering log-linearly in the baseline or in the excitation strength,
all kernels nonnegative --- the MBPP construction is exact. The moment a
\emph{rectified, sign-indefinite} self-interaction is introduced, the MBPP becomes
an \emph{approximation}, and Proposition~\ref{prop:nonlinear} tells us the sign of
the resulting bias. Section~\ref{sec:experiments} deliberately crosses this
boundary, so that we can measure what the approximation costs.

\begin{remark}[lawful anti-excitation]
There \emph{is} a fully lawful way to express ``transactions are suppressed in
some periods'': make the suppression \emph{exogenous}. A covariate
$Z_{\text{cool}}(t)$ (a cooldown / quiet-period indicator) entering the excitation
with a negative coefficient, $\delta_{\text{cool}}<0$, multiplies excitation by
$e^{\delta_{\text{cool}}}<1$ during cooldowns --- anti-excitation that keeps the
intensity nonnegative and the model linear, hence exactly covered by
Theorem~\ref{thm:excitation}. We use both kinds below: lawful exogenous
suppression in the correctly specified experiment, and unlawful endogenous
rectified self-inhibition in the misspecified one.
\end{remark}

% ===========================================================================
\section{The investment case study and experiments}
\label{sec:experiments}

We now test the theory on a synthetic investment market. Every number below is
produced by the companion code (\code{experiments/exp12\_investments.py},
\code{exp13\_grid.py}) and the package estimators; nothing is hand-tuned to
flatter the method.

\subsection{The data-generating family}
\label{sec:dgp}

We simulate $M=3$ sectors. Each sector has a constant immigrant rate; investments
trigger further investments through a kernel $\varphi$; the triggering strength is
modulated by an observable, piecewise-constant \emph{popularity} covariate
$Z(t)$ through $e^{\delta_{\mathrm{pop}}Z(t)}$. We vary three things:
\begin{itemize}[leftmargin=1.4em,itemsep=1pt]
\item the \textbf{kernel}: \emph{poisson} ($\kappa=0$, no self-excitation),
\emph{exp} (single timescale), \emph{sumexp3} (a mixture of three exponential
timescales), \emph{powerlaw} (mild heavy tail $\propto(1+t)^{-4}$), each normalized
to unit mass so the branching ratio is comparable;
\item the \textbf{covariate importance}: $\delta_{\mathrm{pop}}\in\{0,\,0.35,\,0.80\}$
(none / small / large);
\item the \textbf{specification}: we always fit the single-timescale exponential
MBPP, so \emph{exp} is correctly specified and \emph{sumexp3} / \emph{powerlaw} are
mis-specified in kernel shape, while \emph{poisson} tests whether we falsely detect
self-excitation.
\end{itemize}
Data are generated by exact Ogata thinning with a positive-part envelope (valid
because, within a covariate regime and between events, the positive contribution to
the intensity is non-increasing); counts are then interval-censored into $25$--$30$
windows. We fit each sector and report the representative sector, averaged over
i.i.d.\ histories that share the covariate path, with a held-out split.

\subsection{Correct vs.\ mis-specified, across kernels and covariate strength}
\label{sec:grid}

Table~\ref{tab:grid} reports, for every cell, the recovered covariate coefficient
$\hat\delta_{\mathrm{pop}}$, the fitted branching ratio $\hat\kappa$, and the
improvement in held-out interval-censored log-likelihood from adding the covariate
to the excitation (positive $=$ the covariate helps out of sample).
Figure~\ref{fig:grid} visualizes the recovery and the held-out gains.

\begin{table}[t]
\centering\small
\caption{Factorial grid (\code{exp13\_grid.py}). Single-timescale exponential MBPP
fit to each data-generating kernel. $\hat\delta_{\mathrm{pop}}$ is the recovered
covariate effect (truth in the $\delta$ column); ``held-out $\Delta$'' is the
out-of-sample IC-LL improvement from modelling the covariate.}
\label{tab:grid}
\begin{tabular}{@{}llrrrr@{}}
\toprule
kernel & specification & $\delta_{\mathrm{pop}}$ & $\hat\delta_{\mathrm{pop}}$ & $\hat\kappa$ & held-out $\Delta$\\
\midrule
poisson  & correct ($\kappa\!=\!0$) & 0.00 & $+0.01$ & $0.93^{\dagger}$ & $+0.04$\\
poisson  & correct ($\kappa\!=\!0$) & 0.80 & $+0.01$ & $0.93^{\dagger}$ & $+0.04$\\
\midrule
exp      & \textbf{correct}         & 0.00 & $-0.16$ & 0.22 & $+0.12$\\
exp      & \textbf{correct}         & 0.35 & $+0.70$ & 0.23 & $-0.10$\\
exp      & \textbf{correct}         & 0.80 & $+1.05$ & 0.27 & $+1.04$\\
\midrule
sumexp3  & mis-spec.\ (timescales)  & 0.35 & $+0.31$ & 0.31 & $+0.08$\\
sumexp3  & mis-spec.\ (timescales)  & 0.80 & $+0.50$ & 0.44 & $+0.18$\\
\midrule
powerlaw & mis-spec.\ (shape)       & 0.00 & $+0.02$ & 0.91 & $-0.59$\\
powerlaw & mis-spec.\ (shape)       & 0.80 & $+1.78$ & 0.12 & $+0.06$\\
\bottomrule
\end{tabular}
\\[2pt]
{\footnotesize $^{\dagger}$For purely Poisson data with a free constant baseline,
$(\mu,\kappa)$ is unidentified (only the mean rate $\mu/(1-\kappa)$ is constrained),
so $\hat\kappa$ is not meaningful here; the estimator correctly finds no covariate
effect and no held-out gain.}
\end{table}

\begin{figure}[t]
\centering
\includegraphics[width=\linewidth]{../results/exp13_grid.png}
\caption{Grid results: recovered covariate effect (left) and held-out IC-LL
improvement (right) across data-generating kernels and covariate strength.}
\label{fig:grid}
\end{figure}

Three readings, all consistent with the theory.
\emph{(i) Correct specification works.} For the exponential kernel the covariate's
sign is always recovered, the branching ratio sits near its truth
($\hat\kappa\approx0.22$--$0.27$ vs.\ $0.30$), and a strong covariate
($\delta=0.8$) gives a large, unambiguous held-out gain ($+1.04$). The magnitude of
$\hat\delta$ overshoots modestly ($1.05$ vs.\ $0.80$); the overshoot grows under
mis-specification.
\emph{(ii) Mis-specification degrades gracefully, then sharply.} A wrong
\emph{number of timescales} (sumexp3 fit by one exponential) \emph{attenuates}
$\hat\delta$ ($0.50$ vs.\ $0.80$) and inflates $\hat\kappa$ ($0.44$) as the single
exponential strains to mimic the slow tail --- but the covariate still helps out of
sample. A wrong \emph{kernel shape} (power-law) is harder: the exponential fit is
fragile, pushing $\hat\kappa$ toward criticality ($0.9$) in some cells, and
$\hat\delta$ is unreliable. This is the identifiability of Remark~\ref{rem:ridge}
under model error: the covariate is a level effect and survives mild
mis-specification, but a heavy tail the model cannot represent corrupts it.
\emph{(iii) No false positives.} With no covariate ($\delta=0$) the estimate is
$\approx0$; with no self-excitation (poisson) we recover $\hat\delta\approx0$ and no
held-out gain --- though, as the footnote warns, counts with a free constant
baseline cannot by themselves certify $\kappa=0$.

\subsection{A focused study: self-inhibition and lawful anti-excitation}
\label{sec:focused}

The grid keeps self-excitation positive so the per-sector fit is well posed. We now
deliberately cross the boundary of Proposition~\ref{prop:nonlinear}: the
data-generating process has a heavy-tailed power-law kernel, \emph{negative}
self-excitation through a rectified ($\max(0,\cdot)$) intensity, positive
cross-excitation between the similar sectors, and popularity-modulated excitation
(\code{exp12\_investments.py}, regime A). We fit the linear exponential MBPP, with
and without the covariate. As a control (regime B) we generate from a process the
augmented MBPP represents \emph{exactly} --- exponential kernel, covariate-modulated
excitation, plus an exogenous \emph{cooldown} covariate with a negative coefficient
(lawful anti-excitation, Theorem~\ref{thm:excitation}) --- and fit the same
estimator. Table~\ref{tab:focused} summarizes; Figure~\ref{fig:focused} shows both.

\begin{table}[t]
\centering\small
\caption{Focused study (\code{exp12\_investments.py}). Regime A is triply
mis-specified (power-law shape, nonlinear self-inhibition, unmodelled
cross-excitation); regime B is correctly specified.}
\label{tab:focused}
\begin{tabular}{@{}lll@{}}
\toprule
quantity & Regime A (mis-specified) & Regime B (correct)\\
\midrule
covariate $\delta_{\mathrm{pop}}$ & truth $0.70$, $\hat{}=1.39$ & truth $0.60$, $\hat{}=0.675$\\
cooldown $\delta_{\mathrm{cool}}$ & --- & truth $-0.50$, $\hat{}=-0.699$\\
branching $\kappa$ & --- & truth $0.40$, $\hat{}=0.370$\\
decay $\theta$ & --- & truth $1.00$, $\hat{}=0.819$ (weak)\\
baseline $\mu$ & --- & truth $2.20$, $\hat{}=2.233$\\
held-out IC-LL (no-cov $\to$ cov) & $18.88\to18.69$ & ---\\
dispersion & $1.07$ (inhibition tames clustering) & $2.11$ (over-dispersion)\\
\bottomrule
\end{tabular}
\end{table}

\begin{figure}[t]
\centering
\includegraphics[width=0.86\linewidth]{../results/exp12_investments_A.png}\\[4pt]
\includegraphics[width=0.78\linewidth]{../results/exp12_investments_B.png}
\caption{Top (regime A, mis-specified): counts track popularity, the covariate
model improves held-out fit, and $\hat\delta_{\mathrm{pop}}$ has the right sign but
an inflated magnitude. Bottom (regime B, correct): every parameter --- including the
\emph{negative} cooldown coefficient (lawful anti-excitation) --- is recovered.}
\label{fig:focused}
\end{figure}

The control is decisive: when reality matches the model, the augmented estimator
recovers the branching ratio to within $0.03$, the baseline to within $0.03$, and
both covariate coefficients --- crucially including the \emph{negative} cooldown
effect, which is how anti-excitation enters lawfully --- to within $0.10$, with only
the decay $\theta$ left loose, exactly as Remark~\ref{rem:ridge} predicts. Under the
triply mis-specified regime the covariate's sign and significance survive and
modelling it improves out-of-sample likelihood, but its magnitude is inflated
roughly twofold; we attribute this, as in the grid, to the short-memory exponential
being forced to explain the power-law's persistence through the covariate. The
dispersion numbers are a free diagnostic: the rectified self-inhibition makes regime
A \emph{less} over-dispersed than the genuinely branching regime B, a sanity check
that the two regimes differ in the way the model says they should.

\subsection{Genuinely multivariate estimation}
\label{sec:multifit}

The estimator \code{fit\_mbpp\_ic\_excitation\_multi} fits the full
construction of Theorem~\ref{thm:excitation} in $M$ dimensions: an $M$-vector
baseline, an $M\times M$ branching matrix, a shared decay, and a covariate vector
$\delta$, with the compensator from the multivariate LTV solver. Two checks. First,
\emph{exactness of the reduction}: with $\delta=0$ the covariate solver coincides
with the constant-excitation multivariate ODE to $8.9\times10^{-16}$, confirming the
forward model is correct. Second, \emph{recovery}: from interval counts drawn from
the model the estimator recovers the baseline vector and the covariate effect with
the correct sign and the overall branching (spectral radius $\hat{}=0.31$ vs.\
truth $0.30$). The full $M\times M$ \emph{structure} --- which sector excites which
--- is, however, only weakly identified from counts at moderate sample sizes: self-
and cross-excitation both raise a sector's counts, and only cross-interval
correlations disambiguate them. This is the multivariate face of
Remark~\ref{rem:ridge}: levels and the covariate are identified, fine structure
needs more data (or event-time information). The estimator is wired and validated;
publication-scale recovery of the matrix is a matter of sample size and compute, not
of the method.

\subsection{Learned solvers (work in progress)}
\label{sec:learned}

Finally, the same family is the testbed for replacing the numerical solver with a
\emph{learned} one --- a Fourier neural operator or DeepONet for the forward map, a
physics-informed network (PINN/PINO) trained on the MBPP residual as a fast
differentiable surrogate inside the fit, and an amortized network for the inverse
map $\text{counts}\mapsto(\kappa,\theta,\delta)$. The enabling fact is proved in
Section~\ref{sec:augment} and re-confirmed numerically here: the MBPP residual is
\emph{linear in $\xi$}, so its global minimum is the exact solution. We verify that
the residual evaluated at exact solutions is $1.4\times10^{-14}$ and that the
residual's zero reproduces the Volterra solve to $6\times10^{-15}$ --- i.e.\ the
objective a PINN/PINO minimizes is well posed across the whole parameter family.
The full protocol (datasets, models, metrics, success criteria) is specified in
\code{PURE\_LEARNING.md} with runnable scaffolding in \code{experiments/learning/};
those runs require JAX/TensorFlow and are deferred to that pipeline.

% ===========================================================================
\section{Conclusion}
\label{sec:conclusion}

The Mean Behavior Poisson Process turns an intractable problem --- fitting a
self-exciting process to data with no event times --- into a tractable one, by
replacing the random Hawkes intensity with its deterministic mean and invoking
Watanabe's theorem to recover a genuine Poisson likelihood. We have reproduced the
foundational results in full and then established our central methodological point
with the precision a skeptic requires: covariates entering log-linearly in the
baseline preserve the construction \emph{exactly}; covariates entering the
excitation preserve the mean-field identity, the Poisson structure and the
likelihood, trading only the convolution (and the transform machinery that rests on
it) for an exactly solvable linear time-varying ODE; and the single place the
mathematics genuinely breaks --- rectified, sign-indefinite self-inhibition --- is
identified precisely, with the sign of its bias. The experiments bear this out: under
correct specification the augmented estimator recovers branching ratios and covariate
effects (including lawful anti-excitation) cleanly; under mis-specification it
degrades in the way the identifiability analysis predicts, gracefully for timescale
errors and sharply for an unrepresentable heavy tail. The honest limits --- the
weakly identified decay, the weakly identified multivariate fine structure, the
Poisson/Hawkes degeneracy under a free baseline --- are stated as plainly as the
successes.

\paragraph{Acknowledgement of sources.}
Sections~\ref{sec:prelim}--\ref{sec:estim} reproduce, with proofs, the framework of
Rizoiu et al.\ (2022), \emph{Interval-censored Hawkes processes}, JMLR 23(1); the
covariate augmentation of Section~\ref{sec:augment}, the precise nonlinear-inhibition
boundary, and all experiments are the present contribution.

\begin{thebibliography}{9}
\small
\bibitem{ICHawkes2022} M.-A.~Rizoiu, A.~Soen, S.~Li, P.~Calderon, L.~J.~Dong,
A.~K.~Menon, L.~Xie (2022). Interval-censored Hawkes processes. \emph{JMLR}
23(1):1--84.
\bibitem{HawkesOakes1974} A.~G.~Hawkes, D.~Oakes (1974). A cluster process
representation of a self-exciting process. \emph{J.~Appl.~Probab.} 11(3).
\bibitem{Watanabe1964} S.~Watanabe (1964). On discontinuous additive functionals and
L\'evy measures of a Markov process. \emph{Japan.~J.~Math.} 34.
\bibitem{DaleyVereJones} D.~J.~Daley, D.~Vere-Jones (2003). \emph{An Introduction to
the Theory of Point Processes}. Springer.
\bibitem{Banerjee2005} A.~Banerjee, S.~Merugu, I.~Dhillon, J.~Ghosh (2005).
Clustering with Bregman divergences. \emph{JMLR} 6.
\bibitem{Rizoiu2017} M.-A.~Rizoiu et al. (2017). Expecting to be HIP: Hawkes
intensity processes for social media popularity. \emph{WWW}.
\bibitem{GLS1990} G.~Gripenberg, S.-O.~Londen, O.~Staffans (1990). \emph{Volterra
Integral and Functional Equations}. Cambridge.
\bibitem{LiPINO2021} Z.~Li et al. (2021). Physics-Informed Neural Operator for
learning PDEs. \emph{arXiv:2111.03794}.
\end{thebibliography}

\end{document}
